\section{Деревья и отрезки}

\subsection{Домашнее задание}
\begin{enumerate}
  \item
    На плоскости даны $n$ точек и $m$ прямоугольников с целочисленными координатами. Стороны прямоугольников
    параллельны осям координат. Найдите количество точек, покрытых ровно $k$ прямоугольниками. $\O((n + m)\log{m})$.

    \textbf{Решение:}

    \begin{itemize}
      \item Идея: используем метод сканирующей прямой (sweep line) по оси $x$ совместно с деревом отрезков по оси $y$.

      Каждый прямоугольник $(x_1, y_1, x_2, y_2)$ можно представить как два события:
      \begin{itemize}
        \item при $x = x_1$: прибавить $+1$ на отрезке $[y_1, y_2]$ (прямоугольник начался)

        \item при $x = x_2 + 1$: прибавить $-1$ на отрезке $[y_1, y_2]$ (прямоугольник закончился)
      \end{itemize}

      Тогда в любой момент значение в позиции $y$ дерева отрезков будет равно количеству прямоугольников,
      покрывающих текущую $y$-координату при текущем значении $x$.

      \item Для эффективной работы необходимо координатное сжатие по $y$:

      Соберём все $y$-координаты из прямоугольников ($y_1$, $y_2$ и $y_2 + 1$ для каждого), отсортируем и заменим на ранги.
      Всего различных $y$-значений не более $3m$, поэтому размер дерева отрезков равен $\O(m)$.
      
      Для каждой точки $(p_x, p_y)$ найдём её позицию в сжатых координатах бинарным поиском за $\O(\log m)$.

      \item Дерево отрезков используем в режиме range update + point query с lazy propagation:

      \begin{itemize}
        \item Обновление отрезка $[l, r]$ на $+\delta$: добавляем $\delta$ в поле $lazy$ вершины, покрывающей нужный отрезок

        \item Запрос в точке $i$: спускаемся до листа, проталкивая отложенные обновления, и возвращаем значение
      \end{itemize}

      \item Финальный алгоритм:

      \begin{enumerate}

        \item Собрать $y$-координаты из прямоугольников и выполнить координатное сжатие. Построить дерево отрезков размера $\O(m)$.

        \item Для каждого прямоугольника $(x_1, y_1, x_2, y_2)$ создать два события:
        \begin{itemize}
          \item $(x_1,\; +1,\; y_1,\; y_2)$ -- открытие
          \item $(x_2 + 1,\; -1,\; y_1,\; y_2)$ -- закрытие
        \end{itemize}
        Отсортировать $2m$ событий по $x$-координате.

        \item Для каждой точки $(p_x, p_y)$ бинарным поиском по отсортированному массиву событий за $\O(\log m)$
        найдём позицию $p_x$ среди $x$-координат событий. По этой позиции определим, после какого события нужно обработать данную точку
         (то есть после применения всех событий с $x \leq p_x$, но до событий с $x > p_x$). Сгруппируем точки по этим позициям.

        \item Проходим по отсортированным событиям слева направо. После применения очередной группы событий с одинаковой $x$-координатой
        отвечаем на все точки, привязанные к данной позиции: для каждой такой точки делаем point query в дереве отрезков по сжатой $p_y$ и проверяем, 
        равно ли значение $k$.

      \end{enumerate}

      \item Корректность:
      \begin{itemize}
        \item В момент обработки точки $(p_x, p_y)$ все события с $x$-координатой $\leq p_x$ уже применены к дереву отрезков, 
        а события с $x$-координатой $> p_x$ ещё нет.
        
        \item Тогда все прямоугольники с $x_1 \leq p_x$ уже открыты (событие $+1$ применено), а все прямоугольники с $x_2 < p_x$ 
        уже закрыты (событие $-1$ при $x_2 + 1 \leq p_x$ применено). Прямоугольники с $x_1 > p_x$ ещё не открыты.
        
        \item Таким образом, значение в позиции $p_y$ дерева отрезков равно числу прямоугольников, покрывающих точку $(p_x, p_y)$.
      \end{itemize}

      \item Сложность по времени:
      \begin{itemize}
        \item Координатное сжатие по $y$: сортировка $\leq 3m$ значений -- $\O(m \log m)$

        \item Сортировка $2m$ событий по $x$ -- $\O(m \log m)$

        \item Распределение $n$ точек по интервалам (бинарный поиск) -- $\O(n \log m)$

        \item Обработка $2m$ событий (range update на дереве отрезков размера $\O(m)$) -- $\O(m \log m)$

        \item Обработка $n$ запросов (point query на дереве отрезков размера $\O(m)$) -- $\O(n \log m)$
      \end{itemize}

      Итого: $\O(m \log m + n \log m) = \O((n + m) \log m)$

      \item Сложность по памяти:
      \begin{itemize}
        \item Дерево отрезков: $\O(m)$
        \item Массив событий: $\O(m)$
        \item Списки точек по интервалам: $\O(n)$ суммарно
      \end{itemize}

      Итого: $\O(n + m)$

    \end{itemize}

  \item
    Дана последовательность чисел длины $n$. Придумайте как за $\O(n\log{n} + m\sqrt{n}\log{n})$ ответить
    на $m$ следующих запросов:
    \begin{itemize}
      \item \texttt{sum(l, r)} -- посчитать сумму на отрезке $[l,r]$.

      \item \texttt{inc(l, r, x)} -- увеличить все значения отрезка $[l,r]$, меньшие $x$, до $x$.
    \end{itemize}

    \textbf{Решение:}

    \begin{itemize}

      \item Идея: используем корневую декомпозицию с блоками размера $B = \sqrt{n}$.
      Для каждого блока храним отсортированную копию его элементов и массив префиксных сумм отсортированного массива,
      а также отложенное значение $floor$ (операция \texttt{inc} -- это по сути операция $\max$ поэлементно).

      \item Структуры данных для каждого блока $b$:
      \begin{itemize}
        \item $array$ -- исходный массив элементов блока

        \item $sorted$ -- отсортированная копия элементов блока

        \item $prefix[i]$ -- сумма первых $i$ элементов массива $sorted$: $prefix[i] = \sum_{j=0}^{i-1} sorted[j]$

        \item $f$ -- отложенное значение floor (изначально $-\infty$). Означает, что все элементы блока, меньшие $f$, должны быть заменены на $f$.
      \end{itemize}

      \item Вспомогательные операции:

      \begin{itemize}

        \item \texttt{push(b)} -- проталкивание отложенного floor в блок $b$:
        \begin{itemize}
          \item для каждого элемента $array[i]$ блока выполним $array[i] := \max(array[i], f)$

          \item затем пересортируем блок и пересчитаем $prefix$

          \item сбросим $f := -\infty$
        \end{itemize}
        
        Стоимость: $\O(B \log B) = \O(\sqrt{n} \log n)$.

        \item \texttt{block\_sum(b)} -- сумма элементов блока $b$ с учётом floor $f$:
        \begin{itemize}
          \item бинарным поиском в $sorted$ найдём индекс $j$ -- количество элементов, меньших $f$.
          
          \item тогда сумма блока равна $j \cdot f + (prefix[B] - prefix[j])$: первые $j$ элементов заменяются на $f$,
        остальные остаются без изменений
        \end{itemize}

        Стоимость: $\O(\log B) = \O(\log n)$.

      \end{itemize}

      \item Операция \texttt{sum(l, r)}:
      \begin{itemize}

        \item Для крайних (неполных) блоков:
        \begin{itemize}
          \item вызываем $push$
          \item затем суммируем нужные элементы напрямую
        \end{itemize}
        Стоимость: $\O(\sqrt{n} \log n)$ на каждый из $\leq 2$ крайних блоков.

        \item Для полных блоков: вызываем $block\_sum$
        
        Стоимость: $\O(\log n)$ на блок, всего $\O(\sqrt{n})$ полных блоков.
      \end{itemize}
      Итого на запрос: $\O(\sqrt{n} \log n)$.

      \item Операция \texttt{inc(l, r, x)}:
      \begin{itemize}

        \item Для крайних блоков:
        \begin{itemize}
          \item вызываем $push$
          \item затем поэлементно применяем $array[i] := \max(array[i], x)$
          \item после чего пересортируем блок и пересчитаем $prefix$
        \end{itemize}
        
        Стоимость: $\O(\sqrt{n} \log n)$ на каждый из $\leq 2$ крайних блоков.

        \item Для полных блоков: обновляем $f := \max(f, x)$.
        
        Стоимость: $\O(1)$ на блок, всего $\O(\sqrt{n})$ полных блоков.

      \end{itemize}
      Итого на запрос: $\O(\sqrt{n} \log n)$.

      \item Сложность по времени:
      \begin{itemize}
        \item Предподсчёт: сортировка каждого из $n / \sqrt{n} = \sqrt{n}$ блоков размера $\sqrt{n}$: \newline
        $\O(\sqrt{n} \cdot \sqrt{n} \log \sqrt{n}) = \O(n \log n)$

        \item Каждый запрос: $\O(\sqrt{n} \log n)$
      \end{itemize}

      Итого: $\O(n \log n + m\sqrt{n} \log n)$

      \item Сложность по памяти:
      \begin{itemize}
        \item Массив $array$: $\O(n)$

        \item Отсортированные копии и префиксные суммы: $\O(n)$

        \item Значения $f$ для каждого блока: $\O(\sqrt{n})$
      \end{itemize}

      Итого: $\O(n)$

    \end{itemize}

  \item
    Дана последовательность длины $n$. Элементы последовательности --- целые числа, не превосходящие $C$.
    Научитесь обрабатывать online-запросы:

    \begin{itemize}
      \item Циклически сдвинуть подотрезок последовательности на 1 влево или вправо.
        Например, в последовательности (1, 2, 3, 4, 5) циклический сдвиг
        (2, 3, 4) на 1 вправо даст последовательность (1, 4, 2, 3, 5)

        \item Посчитать, сколько раз встречается число $x$ на подотрезке последовательности.
    \end{itemize}
    $\O(\sqrt{n})$ на запрос, предподсчёт за $\O(C\sqrt{n} + n)$.

    \textbf{Решение:}

    \begin{itemize}

      \item Идея: используем корневую декомпозицию с блоками размера $B = \sqrt{n}$.
      Каждый блок хранит элементы в деке, а также массив счётчиков $count[x]$ -- количество вхождений значения $x$ в данном блоке.

      \item Структуры данных для каждого блока $b$:
      \begin{itemize}
        \item $deque$ -- кольцевой буфер размера $B$

        \item $count[0 \ldots C]$ -- массив счётчиков: $count[x]$ = количество элементов блока, равных $x$
      \end{itemize}

      \item Операция \texttt{count(l, r, x)} -- подсчёт вхождений $x$ на $[l, r]$:
      \begin{itemize}
        \item Для крайних (неполных) блоков: перебираем элементы и считаем вхождения $x$ вручную.
        Стоимость: $\O(\sqrt{n})$ на $\leq 2$ крайних блока.

        \item Для полных блоков: прибавляем $count_b[x]$. Стоимость: $\O(1)$ на блок, всего $\O(\sqrt{n})$ полных блоков.
      \end{itemize}
      Итого: $\O(\sqrt{n})$.

      \item Операция \texttt{shift\_right(l, r)} -- циклический сдвиг $[l, r]$ вправо на 1:

      Эквивалентно: извлечь элемент $val = array[r]$ и вставить его в позицию $l$, сдвинув остальные элементы вправо.

      \begin{enumerate}
        \item В блоке, содержащем позицию $r$: извлекаем элемент из позиции $r$ внутри блока (сдвигая элементы правее влево).
        Блок теряет один элемент. Обновляем $count[val] := count[val] - 1$. Стоимость: $\O(\sqrt{n})$.

        \item Каскад по полным блокам (от блока $r$ к блоку $l$): каждый блок отдаёт свой последний элемент следующему блоку справа
        ($pop\_back$), а от блока слева получает элемент в начало ($push\_front$).
        Обновляем $count$ для входящего ($+1$) и исходящего ($-1$) элементов.
        Стоимость: $\O(1)$ на блок, $\O(\sqrt{n})$ блоков суммарно.

        \item В блоке, содержащем позицию $l$: вставляем $val$ в позицию $l$ внутри блока (сдвигая элементы правее вправо).
        Обновляем $count[val] := count[val] + 1$. Стоимость: $\O(\sqrt{n})$.
      \end{enumerate}

      После каскада все блоки сохраняют исходный размер $B$, так как суммарно извлечён и вставлен ровно один элемент.

      Итого на запрос: $\O(\sqrt{n})$.

      \item Операция \texttt{shift\_left(l, r)} -- аналогична, но каскад идёт в обратном направлении:
      извлекаем $array[l]$, передаём элементы через $pop\_front / push\_back$ слева направо, вставляем в позицию $r$.
      Итого: $\O(\sqrt{n})$.

      \item Сложность по времени:
      \begin{itemize}
        \item Предподсчёт: для каждого из $\sqrt{n}$ блоков инициализируем массив $count[0 \ldots C]$ нулями за $\O(C)$
        и заполняем за $\O(\sqrt{n})$.
        Итого: $\O(\sqrt{n} \cdot C + n) = \O(C\sqrt{n} + n)$.
        
        \item Каждый запрос (сдвиг или подсчёт): $\O(\sqrt{n})$.
      \end{itemize}

      \item Сложность по памяти:
      \begin{itemize}
        \item Деки для блоков: $\O(n)$ суммарно

        \item Массивы $count$: $\O(C \cdot \sqrt{n})$
      \end{itemize}

      Итого: $\O(n + C\sqrt{n})$

    \end{itemize}

  \item
    Дана последовательность из $n$ чисел и $m$ отрезков $[l_i, r_i]$ на ней.
    Про каждый отрезок надо определить число инверсий на нем. $\O((m + n) \sqrt{n} \log{n})$.

\end{enumerate}
